\SourceLeaf{378}
\par\bigskip\Needspace{9\baselineskip}
\pdfbookmark[1]{Theorie der biquadratischen Reste. VI.}{nachlass-biquadratic-VI}
\centerline{[VI.]}\medskip
\begin{center}\large\scshape Theorie der biquadratischen Reste.\end{center}
\BookRule{.18\textwidth}
\renewcommand{\Article}[1]{\par\medskip\Needspace{4\baselineskip}\pdfbookmark[2]{Art. #1}{nachlass-biquadratic-VI-art-#1}\centerline{\textbf{#1.}}\nobreak\smallskip}
\Article{1}
Eine Folge komplexer ganzer Zahlen $\varphi,\varphi',\varphi''$ usw. sei so beschaffen, dass sie erstens modulo $\mu=\alpha+\beta i$ alle paarweise inkongruent sind, wobei $\alpha$ und $\beta$ ganze reelle Zahlen bezeichnen, und dass zweitens jede komplexe ganze Zahl modulo $\mu$ einer von ihnen kongruent ist. Unter dieser Voraussetzung heißt die Gesamtheit der Zahlen $\varphi,\varphi',\varphi''$ usw. ein vollständiges Restsystem für den Modul $\mu$. Es ist bewiesen, dass die Anzahl seiner Zahlen gleich der Norm von $\mu$ ist, also der Zahl $\alpha\alpha+\beta\beta$, die mit $\nu$ bezeichnet werden soll.

\Article{2}
Unter den Zahlen $\varphi,\varphi',\varphi''$ usw. ist eine durch $\mu$ teilbar. Wird sie ausgeschlossen und die Gesamtheit der übrigen mit $\chi$ bezeichnet, so bildet $\chi$ ein vollständiges System der durch den Modul nicht teilbaren Reste, deren Anzahl $=\nu-1$ ist. Beschränken wir die Untersuchung auf ungerade Moduln, so ist $\nu-1$ durch $4$ teilbar. Dann sind ferner $f,if,-f,-if$ untereinander inkongruent; folglich sind diejenigen Zahlen in $\chi$, die jeweils $if,-f,-if$ kongruent sind, voneinander und von $f$ verschieden. (Assoziierte und zusammengesetzte Zahlen.)

Daraus ergibt sich eine Zerlegung von $\chi$ in vier Gruppen oder Teilsysteme $\kappa,\kappa',\kappa'',\kappa'''$. Man setzt eine beliebige Zahl aus $\chi$, zum Beispiel $\varphi$, in die Gruppe $\kappa$ und die drei den Zahlen $i\varphi,-\varphi,-i\varphi$ kongruenten Glieder von $\chi$ der Reihe nach in die Gruppen $\kappa',\kappa'',\kappa'''$. Nachdem diese vier Zahlen aus $\chi$ gestrichen sind, setzt man wieder eine beliebige der übrigen in $\kappa$ und die drei auf dieselbe Weise von ihr abhängigen in $\kappa',\kappa'',\kappa'''$. So fährt man fort, bis das ganze System $\chi$ verteilt ist. Die Gruppen $\kappa,\kappa',\kappa'',\kappa'''$ sollen zusammengehörige Viertelsysteme heißen. Offenbar haben sie folgende Eigenschaften:

\SourceLeaf{379}
\begin{enumerate}
\renewcommand{\labelenumi}{\arabic{enumi})}
\item Jedes Viertelsystem besteht aus $\frac14(\nu-1)=\frac14(\alpha\alpha+\beta\beta-1)$ Zahlen.
\item Kennzeichnend für ein Viertelsystem ist, dass keine seiner Zahlen, weder selbst noch mit $i,-1$ oder $-i$ multipliziert, einer anderen Zahl aus demselben Viertelsystem kongruent ist, dass aber jede durch $\mu$ nicht teilbare Zahl, entweder selbst oder mit $i,-1$ oder $-i$ multipliziert, darin vorkommt oder einer seiner Zahlen kongruent ist.
\item So wie die Multiplikation der Zahlen in $\kappa$ mit $i,-1$ und $-i$ jeweils die Zahlen in $\kappa',\kappa'',\kappa'''$ oder ihnen kongruente Zahlen ergibt, erzeugt die Multiplikation der Zahlen in $\kappa'$ mit jenen Faktoren jeweils wieder die Zahlen in $\kappa'',\kappa''',\kappa$; die Multiplikation der Zahlen in $\kappa''$ erzeugt auf ähnliche Weise die Zahlen in $\kappa''',\kappa,\kappa'$; schließlich erzeugt die Multiplikation der Zahlen in $\kappa'''$ wieder $\kappa,\kappa',\kappa''$.
\end{enumerate}

Der Kürze halber lässt sich diese wechselseitige Abhängigkeit der vier Viertelsysteme so ausdrücken:
\[
\begin{aligned}
\kappa'&\equiv i\kappa,\\
\kappa''&\equiv-\kappa\equiv i\kappa',\\
\kappa'''&\equiv i\kappa''\equiv-\kappa'\equiv-i\kappa.
\end{aligned}
\]
